\usepackage[round,authoryear]{natbib}
\usepackage{paralist}
\usepackage{pgfpages}
\usepackage{mathtools}
\usepackage{amsmath}
\usepackage{amssymb}
\usepackage{xspace}        
\usepackage{graphicx}
\usepackage{listings}
\usepackage{bm}

\bibliographystyle{jss}
%\bibliographystyle{apalike}
%\bibliographystyle{dcu}

\usefonttheme[onlymath]{serif}
% use article-like math letters, from http://tex.stackexchange.com/questions/34265/how-to-get-beamer-math-to-look-like-article-math

\newcommand{\defeq}{\vcentcolon=}

%%% notation for ARMA models
\newcommand\ar{\phi}
\newcommand\ma{\psi}
\newcommand\AR{\Phi}
\newcommand\MA{\Psi}

%%% notation for linear Gaussian POMP models
\newcommand\matA{\mathbb{A}}
\newcommand\matB{\mathbb{B}}
\newcommand\matH{\mathbb{H}}
\newcommand\covmatX{\mathbb{U}}
\newcommand\covmatY{\mathbb{V}}

\newcommand\plim{\overset{p}{\rightarrow}}
\newcommand\Plim{\overset{P}{\rightarrow}}
\newcommand\dlim{\overset{d}{\rightarrow}}
\newcommand\Dlim{\overset{D}{\rightarrow}}

\newcommand\prob[1]{\mathbb{P}\left[{#1}\right]}
\newcommand\expect[1]{\mathbb{E}\left[{#1}\right]}
\newcommand\var[1]{\mathrm{Var}\left[{#1}\right]}
\newcommand\myvar{\mathrm{Var}}
\newcommand\Var{\myvar}
\newcommand\dist[2]{\mathrm{#1}\left(#2\right)}
\newcommand\lik{\mathcal{L}}
\newcommand\loglik{\ell}
\newcommand\R{\mathbb{R}}
\newcommand\N{\mathbb{N}}
\newcommand\Rzero{\mathfrak{R}_0}
\newcommand\argmax{\mathop{\mathrm{argmax}}}
\newcommand\argmin{\mathop{\mathrm{argmin}}}

\newcommand\code[1]{\texttt{#1}}
\newcommand\package[1]{\textbf{#1}}
\newcommand\data[1]{#1^*}
%\newcommand\data[1]{#1}
%\newcommand\estimate[1]{\data{#1}}
\newcommand\estimate[1]{#1}
\newcommand\predict[1]{#1^{P}}
\newcommand\residual{r}

\newcommand\link[2]{\href{#1}{#2}}
\newcommand{\doi}[1]{\link{https://doi.org/#1}{\texttt{doi:~{#1}}}}

\newcommand{\scinot}[2]{#1{\times}10^{#2}}
\newcommand{\dd}[1]{\mathrm{d}{#1}}
\newcommand{\pd}[3][]{%
  \def\ord{#1} \ifx\ord\empty%
  \frac{\partial{#2}}{\partial{#3}}%
  \else \frac{\partial^{#1}{#2}}{\partial{#3}^{#1}}%
  \fi
}
\newcommand{\deriv}[3][]{%
  \def\ord{#1} \ifx\ord\empty%
  \frac{\dd{#2}}{\dd{#3}}%
  \else \frac{\dd^{#1}{#2}}{\dd{#3}^{#1}}%
  \fi
}

\newcommand\Rlanguage{\textsf{R}\xspace}

\newcommand\question{{\bf Question}}
\newcommand\mysolution{{\bf Solution}}
\newcounter{Qcounter}
\newcommand\myquestion{{\stepcounter{Qcounter} \bf Question \CHAPTER.\theQcounter}}
\newcounter{Ecounter}
\newcommand\myexercise{{\stepcounter{Ecounter} \bf Exercise \CHAPTER.\theEcounter}}

\newcommand\myexample{{\bf Example}}
\newcommand\mydot{{\,\cdot\,}}
\newcommand\myref[1]{\m{#1}}

\newcommand\Rspace{\mathcal{R}}

\renewcommand\vec[1]{\boldsymbol{\mathrm{#1}}}
\newcommand\vect[1]{\vec{#1}}
\newcommand\mat[1]{\mathbb{#1}}
\newcommand\pr{\mathbb{P}}
\newcommand\E{\mathbb{E}}

\newcommand\profileloglik[1]{\ell^\mathrm{profile}_#1}
\newcommand\Real{\mathbb{R}}

\newcommand\bi{\begin{itemize}}
\newcommand\ei{\end{itemize}}

\newcommand\normal{\mathrm{normal}}

\newcommand\iid{iid}
\newcommand\MVN{\mathrm{MVN}}
\newcommand\SE{\mathrm{SE}}

\newcommand\pval{\mathrm{pval}}
% \newcommand\var{\mathrm{Var}}
\newcommand\sd{\mathrm{SD}}
\newcommand\sdSample{\mathrm{sd}}
\newcommand\varSample{\mathrm{var}}
\newcommand\cov{\mathrm{Cov}}
\newcommand\Cov{\mathrm{Cov}}
\newcommand\covSample{\mathrm{cov}}
\newcommand\corSample{\mathrm{cor}}
\newcommand\cor{\mathrm{Cor}}
\newcommand\Cor{\cor}
\newcommand\given{{\, | \,}}
\newcommand\param{\,;}
\newcommand\params{\param}
\newcommand\equals{{ \, = \, }}
\newcommand\transpose{{\raisebox{0.5mm}{\mbox{\scriptsize \textsc{t}}}}}
%\newcommand\transpose{\scriptsize{T}}
\newcommand\mycolon{{\hspace{0.5mm}:\hspace{0.5mm}}}
%\newcommand\mycolon{\,{:}\,}
\newcommand\myemph[1]{{\textbf{#1}}}
\newcommand\mymathenv[1]{\textcolor{blue}{#1}}


%% following a style where all math is in blue
%% not currently used 
\newcommand\mymath[1]{\begin{math}\textcolor{blue}{#1}\end{math}}
\newcommand\m[1]{\mymath{#1}}
\newcommand\mydisplaymath[1]{\begin{displaymath}\textcolor{blue}{#1}\end{displaymath}}
\newcommand\myeqnarray[1]{\textcolor{blue}{\begin{eqnarray*}#1 \end{eqnarray*}}}
\newcommand\myspace{\quad}
\newcommand\altdisplaymath[1]{\vspace{1mm}\textcolor{blue}{\begin{math}\displaystyle #1 \end{math}}\vspace{1mm}}

%%%%%%%%%%%%%%%%%


\setlength{\parskip}{0mm}
\setlength{\parindent}{0mm}
\lstset{language=C}

\newcommand\negBeforeCode{}
\newcommand\negAfterCode{}

% \mode<beamer>{\usetheme{AnnArbor}}
\usetheme{metropolis}
% \mode<beamer>{\usecolortheme{dolphin}}
% \mode<beamer>{\usefonttheme{serif}}
\mode<beamer>{\setbeamertemplate{footline}}
\mode<beamer>{\setbeamertemplate{footline}[frame number]}
\mode<beamer>{\setbeamertemplate{frametitle continuation}[from second][\insertcontinuationcountroman]}
\mode<beamer>{\setbeamertemplate{navigation symbols}{}}
\mode<beamer>{\renewcommand{\bibsection}{}}

\mode<handout>{\pgfpagesuselayout{2 on 1}[letterpaper,border shrink=5mm]}


% \mode<beamer>{\usetheme{AnnArbor}}
\mode<beamer>{\usetheme{metropolis}}
\mode<beamer>{\metroset{block=fill}}
% \mode<beamer>{\usecolortheme{wolverine}}

\mode<beamer>{\setbeamertemplate{section in toc}[sections numbered]}
\mode<beamer>{\setbeamertemplate{subsection in toc}[subsections numbered indented]}

% \mode<beamer>{\usefonttheme{serif}}
\mode<beamer>{\setbeamertemplate{footline}}
\mode<beamer>{\setbeamertemplate{footline}[frame number]}
\mode<beamer>{\setbeamertemplate{frametitle continuation}[from second][\insertcontinuationcountroman]}
\mode<beamer>{\setbeamertemplate{navigation symbols}{}}

\mode<handout>{\pgfpagesuselayout{2 on 1}[letterpaper,border shrink=5mm]}

\newcommand\CHAPTER{4}
% \newcommand\answer[2]{\textcolor{blue}{#2}} % to show answers
% \newcommand\answer[2]{\textcolor{red}{#2}} % to show answers
 \newcommand\answer[2]{#1} % to show blank space

\title{\vspace{2mm} \link{https://jeswheel.github.io/4450_f25/}{Mathematical Statistics I}\\ \vspace{2mm}
Chapter \CHAPTER: Expected Values}
\author{Jesse Wheeler}
\date{}

\setbeamertemplate{footline}[frame number]




\begin{document}

\maketitle

\mode<article>{\tableofcontents}

\mode<presentation>{
  \begin{frame}{Outline}
    \tableofcontents
  \end{frame}
}

\section{Discrete random variables}

\begin{frame}{Introduction}
  \begin{itemize}
    \item This material comes primarily from \citet[][Chapter~4]{rice07}.
    \item We will cover the ideas of expected value, variance, as well has higher-order moments.
    \item This includes topics such as conditional expectation, which is one of the fundamental ideas behind many branches of statistics and machine learning.
    \item For instance, most regression / prediction algorithms are built with the idea of minimizing some conditional expectation.
  \end{itemize}
\end{frame}

\begin{frame}[allowframebreaks]{Expectation: Discrete random variables}
  % \begin{itemize}
  %   \item We will begin by defining the expectation for discrete random variables.
  % \end{itemize}
  \begin{block}{Definition: Expectation of discrete random variables}
    Let $X$ be a discrete random variable with pmf $p(x)$, which takes values in the space $\mathcal{X}$. The \alert{expected value} of $X$ is 
    % $$
    % E\big(g(X)\big) = \sum_{x\in \mathcal{X}}g(x)\,p(x).
    % $$
    % In particular, for $g(x) = x$, we have
    $$
    E(X) = \sum_{x \in \mathcal{X}}x\, p(x),
    $$
    provided that $\sum_{x \in \mathcal{X}} |x|\, p(x) < \infty$; otherwise, the expectation is not defined.
  \end{block}
  \begin{itemize}
    \item This is not the most mathematically precise definition of expectation, but a more complete treatment of the topic is outside the scope of this course \citep[See][]{resnick19}.
    \item The concept of the expected value parallels the notion of a \emph{weighted average}.
    \item That is, we weight each possibility $x \in \mathcal{X}$ by their corresponding probability: $\sum_x x \, p(x)$.
    \item $E(X)$ is also referred to as the \alert{mean} of $X$, and is typically denoted $\mu$ or $\mu_X$. 
    \item If the function $p$ is thought of as a weight, then $E(X)$ is the center; that is, if we place the mass $p(x_i)$ at the points $x_i$, then the balancing point is $E(X)$.
    \item Like with the pmf and cdf, we often use subscripts to denote which probability law we are using for the expectation, it if is not clear: $E_X(X)$. 
  \end{itemize}
  
  \begin{exampleblock}{Roulette}
    A roulette wheel has the numbers $1$ through $36$, as well as $0$ and $00$. If you bet \$1 that an odd number comes up, you win or lose \$1 according to whether that event occurs. If $X$ denotes your net gain, $X = 1$ with probability $18/38$ and $X = -1$ with probability $20/28$. The expected value of $X$ is
    $$
    E(X) = 1 \times \frac{18}{38} + (-1) \times \frac{20}{38} = -\frac{1}{19}.
    $$
  \end{exampleblock}
  
  \begin{itemize}
    \item As you might imagine, the expected value coincides in the limit with the actual average loss per game, if you play many games (Chapter~5).
    \item Most casino games have a negative expected value by design; you may win some money, but if a large number of games are played, the house will come out on top. 
  \end{itemize}
  
  \framebreak
  
  \begin{exampleblock}{Geometric Random Variable}
   Suppose that items are produced in a plant are independently defective with probability $p$. If items are inspected one by one until a defective item is found, then how many items must be inspected on average?
  
  \mode<article>{
  Let $X$ denote the number of items inspected, up-to and including the first defective item. $X$ is geometrically distributed, which as pmf
  $$
  p(k) = P(X = k) = p\,(1-p)^{k-1}.
  $$
  Therefore
  \begin{align*}
    E(X) &= \sum_{k = 1}^\infty k p\, (1-p)^{k-1} \\
         &= p \sum_{k = 1}^{\infty} k \, (1-p)^{k-1}.
  \end{align*}
  To work out this summation, we will use a trick that is sometimes useful for infinite series.
  First, lets define $q = 1-p$, and note that $0 < q < 1$. Then, the sum becomes
  $$
  E(X) = p\sum_{k = 1}^\infty k\,q^{k-1}.
  $$
  You might notice that the summand is a power-rule derivative:
  $$
  \frac{d}{dq} q^k = k\,q^{k-1}.
  $$
  This fact is going to be useful, because the left-hand side of this derivative equation is a geometric sum, which we know how to calculate:
  $$
  \sum_{k = 1}^\infty q^k = \sum_{k = 1}^\infty q\,q^{k-1} = q\sum_{j = 0}^\infty q^j = \frac{q}{1-q}.
  $$
  Thus, what we would like to do is write 
  $$
  \frac{d}{dq}\Big(\frac{q}{1-q}\Big) = \frac{d}{dq}\Big(\sum_{k = 1}^\infty q^{k}\Big) \overset{?}{=} \sum_{k = 1}^\infty \frac{d}{dq} q^k = \sum_{k = 1}^{\infty} kq^{k-1}.
  $$
  Now we can easily calculate the left-hand side to be $\frac{1}{(1-q)^2}$, and therefore we want to make the conclusion
  $$
  \sum_{k = 1}^{\infty} k\, q^{k-1} \overset{?}{=} \frac{d}{dq}\Big(\frac{q}{1-q}\Big) = \frac{1}{(1 - q)^2}.
  $$
  The question is: \textbf{Can we move the derivative inside of the infinite sum?} For this particular case, the answer is \emph{yes}.
  For more details, see Slide~21 from Chapter~3. 
  In this class, all of the sums (and integrals) we will consider will be ``well-behaved" and will satisfy the necessary conditions.
  
  With this sorted out, we can now use our trick to finish the calculation:
  \begin{align*}
    E(X) &= p \sum_{k = 1}^{\infty} k \, (q)^{k-1} \\
         &= p\frac{1}{(1 - q)^2} \\
         &= \frac{p}{p^2} = \frac{1}{p}.
  \end{align*}
  }
  \end{exampleblock}
  
  \mode<presentation>{
  \emph{Solution:}
  }
  
  \framebreak
  
  \begin{exampleblock}{Poisson Distribution}
    The Poisson$(\lambda)$ distribution has pmf $p(k) = \frac{\lambda^k}{k!}e^{-\lambda}$, for all $k \geq 0$. Thus, if $X \sim \text{Pois}(\lambda)$, then what is $E[X]$? \mode<article>{
  \begin{align*}
    E[X] &= \sum_{k = 0}^\infty \frac{k\,\lambda^k}{k!}e^{-\lambda}\\
    &= e^{-\lambda} \sum_{k = 0}^\infty \frac{k\, \lambda^{k-1}\cdot \lambda}{k!} \\
    &= \lambda e^{-\lambda} \sum_{k = 1}^\infty \frac{\lambda^{k-1}}{(k-1)!} \\
    &= \lambda e^{-\lambda} \sum_{j = 0}^{\infty} \frac{\lambda^j}{j!} \\
    &= \lambda e^{-\lambda} e^{\lambda} = \lambda.
  \end{align*}
  }
  \end{exampleblock}
  
  \mode<presentation>{
  \emph{Solution:}
  }
  
  \end{frame}

\section{Continuous random variables}

\begin{frame}[allowframebreaks]{Expectation: Continuous random variables}

  \begin{block}{Definition: Expectation of continuous random variables}
    Let $X$ be a continuous random variable with pdf $f(x)$, which takes values in the space $\mathcal{X}$. The \alert{expected value} of $X$ is 
    $$
    E(X) = \int_{x\in \mathcal{X}}xf(x)\, dx.
    $$
    % In particular, for $g(x) = x$, we have
    % $$
    % E(X) = \sum_{x \in \mathcal{X}}x\, p(x).
    % $$
    provided that $\int_{x\in \mathcal{X}}|x|\, f(x)\, dx < \infty$, otherwise the expectation is undefined. 
  \end{block}

 \begin{itemize}
    \item As before, this is not the most mathematically precise definition of expectation, but a more complete treatment of the topic is outside the scope of this course \citep[See][]{resnick19}.
    \item We can still think of $E(X)$ as the center of mass of the density.
  \end{itemize}

\framebreak

\begin{exampleblock}{Exponential$(\lambda)$ expectation}
  Let $X$ have an Exponential$(\lambda)$ density, with $\lambda > 0$. Thus, the pdf of $X$ is given by
  $$
  f_X(x) = \frac{1}{\lambda}e^{-x/\lambda}, \quad 0 \leq x < \infty
  $$
  Find $E[X]$.
  \mode<article>{
  \vspace{2mm}
  \emph{Solution.}
  \begin{align*}
    E[X] &= \int_x x\, f_X(x)\, dx \\
         &= \int_0^\infty \frac{1}{\lambda} x e^{-x/\lambda}\, dx 
  \end{align*}
  \begin{itemize}
    \item To solve this integral, we can use integration by parts:
    $$
    \int u\, dv = uv - \int v\, du.
    $$
    \item We let $u = x$, and $dv = \frac{1}{\lambda}e^{-x/\lambda}$.
    \item Then, $du = dx$, $v = \int \frac{1}{\lambda}e^{-x/\lambda}\, dx = -e^{-x/\lambda}$.
    \item Plugging this in, we get:
    \begin{align*}
      E[X] &= -x e^{-x/\lambda}\big|_0^\infty - \int^\infty_0 -e^{-x/\lambda}\, dx \\
      &= \int^\infty_0 -e^{-x/\lambda}\, dx \\
      &= \lambda e^{-x/\lambda}\big|_0^\infty = \lambda.
    \end{align*}
  \end{itemize}
  }
\end{exampleblock}

\framebreak

\mode<presentation>{
  \framebreak
  \emph{Solution.}
}

\framebreak

\begin{exampleblock}{Gamma Density}
  If $X$ follows a gamma density with parameters $\alpha$ and $\lambda$, then the pdf of $X$ is
  $$
  f(x) = \frac{\lambda^\alpha}{\Gamma(\alpha)} x^{\alpha - 1}e^{-\lambda x}, \quad x \geq 0.
  $$
  Find $E(X)$.

  \mode<article>{
  \vspace{2mm}
  \emph{Solution:} By definition, the expected value of $X$ is
  $$
  E(X) = \int_{0}^\infty (x) \, \frac{\lambda^\alpha}{\Gamma(\alpha)} x^{\alpha - 1}e^{-\lambda x} \, dx.
  $$
  Combining the factors of $x$ in the integrand, we obtain
  $$
  E(X) = \int_{0}^\infty \frac{\lambda^\alpha}{\Gamma(\alpha)} x^{\alpha}e^{-\lambda x}\,dx.
  $$
  Now we will apply the ``integration by density function" trick: we will re-write the integrand so that it corresponds to the density function of some random variable, and use the fact that the density function must integrate to one. Specifically, note that if we let $\alpha^* = \alpha + 1$, then $\alpha = \alpha^* - 1$, and we can express the integral as:
  \begin{align*}
    E(X) &= \int_{0}^\infty \frac{\lambda^\alpha}{\Gamma(\alpha)} x^{\alpha}e^{-\lambda x}\,dx \\
         &= \int_{0}^\infty \frac{\lambda^\alpha}{\Gamma(\alpha)} x^{\alpha^* - 1}e^{-\lambda x}\,dx \\
         &= \Big(\frac{\lambda^\alpha}{\Gamma(\alpha)}\Big)\Big(\frac{\Gamma(\alpha^*)}{\lambda^{\alpha^*}}\Big)\int_{0}^\infty \frac{\lambda^{\alpha^*}}{\Gamma(\alpha^*)} x^{\alpha* - 1}e^{-\lambda x}\,dx \\
         &= \Big(\frac{\lambda^\alpha}{\Gamma(\alpha)}\Big)\Big(\frac{\Gamma(\alpha^*)}{\lambda^{\alpha^*}}\Big)
  \end{align*}
  Where the last step is a result of the fact that the integrand (and support of the integral) matches the density of a Gamma$(\alpha^*, \lambda)$ distribution. Now using the fact that $\alpha* = \alpha + 1$, and that $\Gamma(x+1) = x\Gamma(x)$, we obtain
  \begin{align*}
    E(X) &= \frac{\lambda^\alpha\, \Gamma(\alpha + 1)}{\Gamma(\alpha)\,\lambda^{\alpha + 1}} \\
         &= \frac{\alpha}{\lambda}
  \end{align*}
  }
\end{exampleblock}

\mode<presentation>{
  \framebreak
  \emph{Solution.}
}

% \framebreak

% \begin{exampleblock}{Normal Distribution}
%   Let $X$ follow a normal distribution with mean $\mu$ and variance $\sigma^2$. Using our definition of expectation, show that $E[X] = \mu$.
%   
%   \mode<article>{
%   \vspace{2mm}
%   \emph{Solution:} Recall that the density of a random variable $X \sim N(\mu, \sigma^2)$ is given by
%   $$
%   f(x) = \frac{1}{\sigma\sqrt{2\pi}}e^{-\frac{1}{2}(x - \mu)^2}, \quad -\infty < x < \infty.
%   $$
%   Thus, the expected value is
%   $$
%   E(X) = \frac{1}{\sigma\sqrt{2\pi}} \int_{-\infty}^\infty x e^{-\frac{1}{2}(x - \mu)^2}\, dx.
%   $$
%   Using a change of variables $z = x - \mu$, we have $dx = dz$, and $x = z + \mu$, 
%   }
%   
% \end{exampleblock}

\end{frame}

\section{Expectation of functions of random variables}

\begin{frame}[allowframebreaks]{Functions of random variables}
  \begin{itemize}
    \item We are often interested in functions of random variables: $Y = g(X)$.
    \item Ideas that we have already covered enable us to calculate $E(Y)$.
    \item For instance, you could use the change-of-variables theorem to get the density of $Y$, then use the definition to calculate $E[Y]$.
    \item Fortunately, we don't have to do this. We can instead calculate $E[Y]$ by integrating (or summing) with respect to $X$:
    $$
    E[g(X)] = \int_{x \in \mathcal{X}} g(x)f(x) \, dx.
    $$
    \item We will justify this for the discrete case.
  \end{itemize}
  
  \framebreak
  
  \begin{block}{Theorem 4.1: Expectation of transformed random variables}
    Suppose that $X$ is a random variable and that $Y = g(X)$ for some function $g$. Then,
    \begin{itemize}
      \item If $X$ is discrete with pmf $p(x)$:
      $$
      E(Y) = \sum_x g(x)\, p(x),
      $$
      provided that $\sum_x |g(x)| p(x) < \infty$.
      \item If $X$ is continuous with pdf $f(x)$:
      $$
      E(Y) = \int_{-\infty}^\infty g(x) f(x)\, dx,
      $$
      provided that $\int |g(x)| f(x) \, dx < \infty$.
    \end{itemize}
  \end{block}
  
  \end{frame}
  
  \begin{frame}[allowframebreaks]{Functions of random variables: proof}
  
  \mode<presentation>{
    \emph{Proof:}
  }
  
  \mode<article>{
    \vspace{2mm}
    \emph{Proof:} By definition of expectation,
    $$
    E(Y) = \sum_i y_i p_Y(y_i).
    $$
    Now let $A_i$ denote the set of $x$'s that are mapped to $y_i$ by $g$. That is, $A_i$ is the pre-image of $y_i$, meaning that $x \in A_i$ if $g(x) = y_i$. Then,
    $$
    p_Y(y_i) = \sum_{x \in A_i} p(x),
    $$
    and we can express the expectation as
    \begin{align*}
      E(Y) &= \sum_i y_i p_Y(y_i) \\
           &= \sum_i y_i \sum_{x \in A_i} p(x) \\
           &= \sum_i \sum_{x \in A_i} y_i\, p(x) \\
           &= \sum_i \sum_{x \in A_i} g(x)\,p(x)\\
           &= \sum_x g(x)\, p(x)
    \end{align*}
    Here, the second to last step is because for all $x \in A_i$, $g(x) = y_i$ by definition. The final step is a result of the fact that the $A_i$ are disjoint, and every $x$ belongs to some $A_i$, and thus the sum over $i$ and $x \in A_i$ is the sum of all $x$.
  }
  
  \framebreak
  
  \begin{itemize}
    \item The proof for the continuous case is similar, but does require a measure-theoretic approach to integration.
    \item One important thing to note is that $g\big(E(X)\big)$ is not usually equal to $E\big(g(x)\big)$.
    \item For example, let $Z$ be a standard normal. We know that $E[Z] = 0$, because it's symmetric. However, $P\big(|Z| > 0) = 1$, thus we can readily deduce that $E\big[|Z|\big] \geq 0 = \big|E[Z]\big|$.
    \item This idea can be extended to show that if for all non-negative random variables $X$ that have finite expectation, if $g(x) \leq x$ for some function $g$, then $E[g(X)] \leq E[X]$.
  \end{itemize}
  
\end{frame}

\begin{frame}[allowframebreaks]{Expected value of indicator functions}
  \begin{itemize}
    \item Another important example of expectations is \alert{indicator} random variables.
    \item For example, suppose that $X$ is a random variable. Then $Y = 1[X \in A]$ for some $A \subset \mathcal{X}$ is a random variable.
    
    \begin{exampleblock}{Indicator Random Variable}
      Let $X$ follow a standard normal distribution, and $A = [-1, 1]$. Then $Y = 1[X \in A]$ is defined as the random variables such that $Y(\omega) = 1$ if $X(\omega) \in A$, and $Y(\omega) = 0$ otherwise.
    \end{exampleblock}
    
    \framebreak
    
    \item Expectations of indicator variables are \alert{probabilities}. Let $Y = 1[X \in A]$. 
    \begin{align*}
    E(Y) &= E\big(1[X \in A]\big) \\
    &= \int_{x \in \mathcal{X}} 1[X \in A]\, f(x) \, dx \\
    & = \int_{x \in A} f(x) \, dx = P(X \in A).
    \end{align*}
    \item This fact is useful for deriving some important inequalities.
    \item First, we will show that the expectations of interest actually exist.
    
    \framebreak
    
    \item Let $X$ be a continuous random variable with expectation $E(X)$. From our definition, this implies that $\int |x|\, f(x)\,dx < \infty$.
    \item Now suppose that for some random variable $Y = g(X)$ such that $|Y| \leq |X|$. Then we can deduce that $\int |y|\, f(x)\,dx < \infty$, and therefore $E[Y]$ exists.
    \item Now suppose that $\varphi$ is a non-decreasing, non-negative function, and that for some $a \in \R$, $\varphi(a) > 0$. Then, for all $x \geq a$, $\varphi(x) / \varphi(a) \geq 1$.
    
    \framebreak
    
    \item Define $Y = 1[X \geq a]$. Note that for all possible outcomes $\omega \in \Omega$, 
    $$
    Y = 1[X \geq a] \leq \varphi(X) / \varphi(a) 1[X \geq a] \leq \varphi(X) / \varphi(a).
    $$
    \item Taking expectations of everything (which we argued preserves inequalities),
    $$
    E\big(1[X \geq a]\big) = P(X \geq a) \leq \frac{E\big[\varphi(X)\big]}{\varphi(a)} = E\left[\varphi(X) / \varphi(a)\right].
    $$
    \item This inequality is known as \alert{Markov's (general) inequality}, and is very useful for bounding the probability of particular events. 
    \item Specifically, if $\varphi(x) = |x|^p$, with $p > 0$, then because $|X|$ is always positive, $\varphi$ is non-negative, non-decreasing, and therefore
    $$
    P(|X| \geq a) \leq \frac{E\big[|X|^p\big]}{a^p},
    $$
    \item If we restrict ourselves to the case where $X$ is non-negative, we get the most standard version of the inequality:
    $$
    P(X \geq a) \leq E(X) / a.
    $$
  \end{itemize}

  \framebreak

  \begin{exampleblock}{Markov's Inequality in Action}
    Suppose that an individual is taken randomly from a population that has an average salary of \$50,000. If we assume that salary from the population is approximately independently and identically distributed, we can provide an upper-bound for the probability that the individual is wealthy.
    
    Let $X_i$ be the salary of individual $i$, randomly drawn from said population. Even though all we know is the average salary, Markov's inequality tells use that:
    $$
    P(X \geq 200,000) \leq \frac{50,000}{200,000} = \frac{1}{4}. 
    $$
  \end{exampleblock}

  \framebreak

  \begin{itemize}
    \item Returning to expectations of functions of random variables, we can extend to the multi-variate case
  \end{itemize}
  
  \begin{block}{Theorem 4.2: functions of multiple variables}
    Suppose that $X_1, \ldots, X_n$ are jointly distributed RVs and $Y = g(X_1, \ldots, X_n)$. Then
    \begin{itemize}
      \item IF $X_i$ are discrete with pmf $p(x_1, \ldots, x_n)$, then
      $$
      E(Y) = \sum_{x_1, \ldots, x_n} g(x_1, \ldots, x_n)p(x_1, \ldots, x_n).
      $$
      \item If $X_i$ are continuous with pdf $f(x_1, \ldots, x_n)$, then
      $$
      E(Y) = \int_{\mathcal{X}_1, \ldots, \mathcal{X}_n}g(x_1, \ldots, x_n)f(x_1, \ldots, x_n)\, dx_1\,\ldots,dx_n.
      $$
    \end{itemize}
    In both cases, we need the sum (or integral) of $|g|$ to converge.
  \end{block}

\framebreak

\begin{itemize}
  \item The proof for the discrete case of Theorem 4.2 follows directly that of Theorem 4.1
  \item An immediate consequence of Theorem~4.2 is the following
\end{itemize}

\begin{block}{Corollary~4.2.1}
  If $X$ and $Y$ are independent random variables, and $g$ and $h$ are fixed functions, then
  $$
  E\big[g(X)h(Y)\big] = \Big(E\big[g(X)\big]E\big[h(Y)\big]\Big),
  $$
  provided that the expectations on the right-hand side exist.
\end{block}

\begin{exampleblock}{Example: Breaking sticks}
  A stick of unit-length is broken randomly (uniformly) in two places. What is the average length of the middle piece?
  
  We will interpret this problem to mean that the locations of the two break-points are independent uniform random variables, $U_1$ and $U_2$, and we need to computing $E|U_1 - U_2|$.
  
\mode<article>{
  \emph{Solution:} Theorem $4.2$ implies that we do not need to find the density function of $U_1 - U_2$. Instead, we just need to integrate $|u_1 - u_2|$ against the joint density: $f(u_1, u_2) = 1$, with $0 \leq u_1, u_2 \leq 1$. Thus
  \begin{align*}
    E|U_1 - U_2| &= \int_{0}^1 \int_{0}^1 |u_1 - u_2| \, du_1\,du_2
  \end{align*}
  Splitting this integral into two regions, one where $u_1 \geq u_2$ and one where $u_2 > u_1$, we get
    \begin{align*}
    E|U_1 - U_2| &= \int_{0}^1 \int_{0}^{u_1} (u_1 - u_2) \, du_2\,du_1 + \int_{0}^1 \int_{u_1}^{1} (u_2 - u_1) \, du_2\,du_1 \\
    &= \frac{1}{3}
  \end{align*}
  Logically, this result makes sense as it suggests that if we have two break points (three pieces) and the break points are uniform and random, the middle piece, on average, will be $1/3$ of the length of the original stick. 
}
\end{exampleblock}

\mode<presentation>{
  \emph{Solution:}
}

\end{frame}

\begin{frame}[allowframebreaks]{Linear Combinations of Random Variables}

\begin{itemize}
  \item A useful property of expectation is that it is a \alert{linear operator}. 
  % \item Specifically, we say that an operator $A$ that maps vector space $U$ to $V$, $A:U\rightarrow V$ is \emph{linear} if 
  % $$
  % A(\alpha x + \beta y) = \alpha A x + \beta A y, \quad x, y \in U, \quad \alpha, \beta \in \R.
  % $$
\end{itemize}

\begin{block}{Theorem 4.3: Linear combinations}
  If $X_1, \ldots, X_n$ are jointly distributed random variables with expectations $E(X_i)$, respectively, and $Y = a + \sum_{i = 1}^n b_i X_i$, then,
  $$
  E(Y) = a + \sum_{i = 1}^n b_i E(X_i).
  $$
\end{block}

\framebreak

\emph{Proof.} \mode<article>{We will show this for the continuous case with $n = 2$. The proof for the discrete case is similar, and this argument is readily extended to the case that $n > 2$. First, we will argue that the expectation is well-defined. By definition,
$$
E|Y| = \int |a + b_1 x_1 + b_2 x_2| f(x_1, x_2)\, dx_1 \, dx_2
$$
and by the triangle inequality $|a + b_1 x_1 + b_2 x_2| \leq |a| + |b_1||x_1| + |b_2||x_2|$, and the fact that $E[X_i]$ exists (which implies that $E|X_i| < \infty$), we see that $E|Y| < \infty$. Now we can calculate the expected value. By Theorem 4.2,
\begin{align*}
  E(Y) &= \int (a + b_1 x_1 + b_2 x_2)f(x_1, x_2)\, dx_1\,dx_2 \\
       &= a \int f(x_1, x_2)\, dx_1\,dx_2 + b_1 \int x_1 f(x_1, x_2)\, dx_1\,dx_2 \\
       & \quad + b_2 \int x_2 f(x_1, x_2)\, dx_1\,dx_2 \\
\end{align*}
Note that the first integral is equal to $1$, because it's the integral of a joint pdf over the support. We'll focus now on the second integral, which is evaluated in a similar way as the third integral due to symmetry.
\begin{align*}
  b_1 \int x_1 f(x_1, x_2)\, dx_1\,dx_2 &= b_1 \int x_1 \left(\int f(x_1, x_2)\, dx_2\right)dx_1 \\
  &= b_1 \int x_1 f(x_1) \, dx_1\\
  &= b_1 E[X_1].
\end{align*}
Thus, applying the same idea to the third integral, we get
$$
E(Y) = a + b_1E(X_1) + b_2E(X_2).
$$
}

\framebreak

\begin{itemize}
  \item The previous theorem is extremely useful for calculating expected values. 
  \item An obvious example is \alert{sums} of random variables, such as the arithmetic average.
  \item It's also useful because some distributions can be expressed as the sum of other distributions.
  \item For instance, we saw in a previous example that the sum of two exponential random variables has a Gamma distribution. Thus, if we know the mean of an exponential, we can readily calculate the mean of a Gamma distribution. 
\end{itemize}

\framebreak

\begin{exampleblock}{Expectation of a binomial distribution}
  Let $Y$ follow a Binomial$(p, q)$ distribution. Find the expected value of $Y$. 
  
  \mode<article>{\emph{Solution}. The pmf of a binomial distribution is given by
  $$
  p(k) = \binom{n}{k} p^k (1-p)^{n-k}.
  $$
  Therefore, to find the expected value directly, we need to calculate the sum
  $$
  E(Y) = \sum_{k = 0}^n k \binom{n}{k} p^k (1-p)^{n-k}.
  $$
  It's not immediately clear how one can calculate this sum directly. Instead, we can use the fact that a binomial random variable is defined by the sum of independent Bernoulli distributed random variables. That is, let $X_1, X_2, \ldots, X_n$ be Bernoulli random variables with parameter $p$. Then
  $$
  Y \overset{d}{=} \sum_{i = 1}^n X_i
  $$
  where the symbol $\overset{d}{=}$ is used to indicate that $Y$ and $\sum_i X_i$ have the same distribution\footnote{they are not necessarily equal to each-other, as they are not necessarily defined using the same probability space. Instead, we only require that they have the same distribution.}. Then it is very easy to calculate the expected value of $Y$, as it is the sum of expected values of $X_i$:
  $$
  E[X_i] = p (1) + (1 - p) (0) = p.
  $$
  Thus, 
  $$
  E[Y] = \sum_i E[X_i] = \sum_i p = np. 
  $$
  }
\end{exampleblock}

\mode<presentation>{
\emph{Solution:}
}

\framebreak

% \begin{itemize}
%   \item An application of the expectation of a binomial distribution is ``shotgun sequencing" in genomics.
%   \item This is a method of trying to figure out the sequence of letters that make up a long segment of DNA.
%   \item It is technically difficult / expensive to sequence the entire segment at once if it is very long.
%   \item Idea: Chop the DNA randomly into many small fragments, sequence each fragment, and then assemble the fragments into one long \alert{contig}.
%   \item If there are many fragments, the hope is that the overlaps can be used to assemble the contig.
% 
% \framebreak
%   \item Suppose that the length of the DNA sequence is $G > 0$, and that there are $N$ fragments, each of length $0 < L < G$.
%   \item For example, the length $G$ might be as large as 100,000, and $L$ about 500.
%   \item Now assume that the left end of each fragment is equally likely to be at positions $1, 2, \ldots, G - L + 1$.
%   \item Some questions that we might as are:
%   \begin{itemize}
%     \item What is the probability that a particular location $x \in \{L, L+1, \ldots, G\}$ is covered by at least one fragment?
%     \item How many fragments are expected to cover a particular location?
%   \end{itemize}
%   \item Note: the positions $\{1, 2, \ldots, L - 1\}$ are left out because they behave slightly differently. Only the fragment that covers position $1$ has its left end position at $1$.
%   \item TODO: finish this example.
% \end{itemize}

\begin{exampleblock}{Example: Baseball Card Collection}
  Suppose that you collect baseball cards, that there are $n$ distinct cards, and that on each trial you are equally likely to get a card of any of the types.
  How many trials would you expect to go through until you had a complete set of cards?
  
  \mode<article>{
  \begin{itemize}
    \item Let $X_1$ denote the number of trials up to and including the trial on which the first type of card is collected. Since our first draw is guaranteed to give us a new type of card, we have $X_1 = 1$.
    \item Now let $X_2$ be the number of trials form that point up to and including the trail on which the next card of a new type is obtained.
    \item We can continue this definition, letting $X_i$ be the number of trials needed to obtain the $i$th type of card, after $i-1$ types of cards have already been obtained. Thus, the total number of trials needed until all types of cards have been obtained is $X = \sum_{i = 1}^n X_i$.
    \item What is the distribution of $X_r$, for $1 \leq r \leq n$?
    \item When counting the number of trials to find the $r$th type, we have already found $r-1$ unique types, leaving $n - r + 1$ types that we have not yet collected. We can see then that $X_r$ is a geometric distributed random variable with $p = (n - r + 1) / n$ representing the probability of success. The expected value of a geometric random variable is $1/p$, and therefore we have
        \begin{align*}
      E[X] &= \sum_{i = 1}^n E[X_i] \\
           &= \sum_{i = 1}^n \frac{n}{n-i+1} \\
           &= n\sum_{i = 1}^n \frac{1}{n-i+1} = n\Big(\frac{1}{n} + \frac{1}{n-1} + \ldots + \frac{1}{1}\Big) \\
           &= n\sum_{k = 1}^n \frac{1}{k}
    \end{align*}
    \item  As far as I am aware, there is not a way to express this final sum succinctly, but we can easily calculate the sum on a computer for moderate values of $n$. For very large values of $n$, there are some additional approximations that can be useful. 
    \item For instance, if we suppose that there are roughly 1200 MLB players (40 players per team, 30 teams), and assume that all players are equally as likely to appear on a card (not a good assumption), then the expected number of cards we would need to purchase until we had a card for every player is $1200 \sum_{k = 1}^{1200} \frac{1}{k}$. In R, we could calculate this using \begin{center} \code{1200 * sum(1/(1:1200))}$=9201.25$ \end{center}
    \item For most $n < 1\times 10^8$, modern computers can calculate this sum almost instantly.
    \item For very large $n$, we could use the famous approximation
    $$
    \sum_{k = 1}^n \frac{1}{k} = \log\,n + \gamma + \epsilon_n,
    $$
    where $\gamma \approx 0.577$ is ``Euler's Constant" (which is often defined as the limit of difference between harmonic series and $\log\, n$), and $\epsilon_n \rightarrow 0$. 
  \end{itemize}
  }
  
\end{exampleblock}


\framebreak

\begin{exampleblock}{Example: Group Testing}
  Suppose that a large number, $n$ of blood samples are screened for a rare disease. If each sample is taken individually, $n$ tests will be required. An alternative approach is group individuals into $m$ groups of size $k$, pool the blood samples for each group together and perform a test on the pooled sample. If the pooled test is negative, we know all individuals in the group do not have the rare disease; however, if the test is positive, we can then do tests on each individual in the smaller group. What is the expected number of tests that will be conducted using this approach?
  
  \mode<article>{
  \begin{itemize}
    \item To solve a problem like this, it's important to give names to quantities of interest.
    \item We have $n$ individuals we need to test, and $m$ groups of size $k$, such that $n = mk$.
    \item Let $X_i$ denote the number of tests conducted on the $i$th group. Thus, the total number of tests is $X = \sum_i X_i$. 
    \item If a group tests negative, then $X_i = 1$. If a group tests positive, then we test all members of the group, so $X_i = k$.
    \item Let's let $p$ denote the probability that an individual tests negative (assuming independence, we could let $p$ be 1 minus the proportion of individuals that have the rare disease).
    \item The probability that a group tests negative is therefore $p^k$; in this case, the total number of tests is 1. The probability that a group tests positive is $1 - p^k$, and in this case, the total number of tests is $k+1$ (one for the group, $k$ for each individual test). 
    \item Thus, the expected number of tests is
    $$
    E[X_i] = p^k + (k + 1) (1 - p^k) = k + 1 - kp^k.
    $$
    \item This expectation is the same for all groups, thus
    $$
    E[X] = \sum_i E[X_i] = mE[X_i] = mk + m - mkp^k = n\Big(1 + \frac{1}{k} - p^K\Big).
    $$
    \item We can see now that the expectation of this group testing scenario is $n$ times a proportion $\Big(1 + \frac{1}{k} - p^K\Big)$. Specifically, if we fix $p$ at 1 minus the rate of disease occurrence in a large population, then the number of tests is a function of group size $k$.
    \item Consider using Desmos to play with the value $p$ (start with something like 0.99) as a function of $k$.
  \end{itemize}
  }
  
  % TODO: Example 4.1.2.D?
  
\end{exampleblock}

\framebreak 

\begin{exampleblock}{Example: Counting DNA ``words"}
  Within DNA patterns, we might be interested in finding the number of times a particular combination of letters (or ``word") occurs in a DNA sequence. This can be useful for determining if a region of DNA has unusually large occurrences of specific sequences. Assume each sequence is randomly composed of letters $A, C, G, T$, and that for each location in the sequence, each letter has probability $1/4$. For example, consider occurrence of the ``word" $TATA$.
  $$
  ACTATATAGATATA
  $$
  In the above sequence, we would count $TATA$ 3 times (counting overlaps). In a sequence of length $N$, what is the expected number of times a word of length $q$ occurs?
  
  \mode<article>{\emph{Solution.} To solve this problem, we will use indicator functions. 
  \begin{itemize}
    \item Let $I_n$ denote the event that the start of a word starts at position $n$ of the sequence, for $n \in \{1, 2, \ldots, N - q + 1\}$. 
    \item Thus, $I_n = 1$ if the word occurs in position $n$, and $I_n = 0$ otherwise. 
    \item The total number of words in a sequence of length $N$ is therefore
  $$
  W = \sum_{n = 1}^{N - q + 1} I_n,
  $$
    \item because $I_n$ only has two possibilities, it is Bernoulli distributed.
    \item Our task is now to find the value of $p$, the probability that a word occurs at position $n$. 
    \item Our assumption that we made is that all letters are independent, and equally as likely to occur at any given position. 
    \item Therefore, the probability that the first letter occurs at position $n$ is $1/4$, and the probability that the second letter occurs at the position $n+1$ is $1/4$, and so on.
    \item By the multiplication principle, $P(I_n = 1) = (1/4)^q$, and the expected value is $E[I_n] = (1/4)^q$.
    \item Thus, the expected value of $W$ is
    $$
    E[W] = \sum_{n = 1}^{N - q + 1} E[I_n] = (N - q + 1)(1/4)^q.
    $$
  \end{itemize}
  }
\end{exampleblock}

\end{frame}

\begin{frame}[allowframebreaks]{Expected value as a predictor}

\begin{itemize}
  \item One useful property of the expectation is that it serves as a good predictor for the value of a random variable. 
  \item Suppose $X$ is a random variable with well-defined expectation, and that we want to make a prediction for the value of $X$.
  \item Denote our predicted value of $X$ as $b$. 
  \item One common way to measure accuracy using the Mean-Squared Error (MSE), which is defined as:
  $$
  \text{MSE}(b) = E\big[(X - b)^2\big].
  $$
  \item Here, the closer $b$ is to $X$, the smaller $(X- b)^2$ is. We take the expectation because $X$ is random.
  \item By this measure, the best predictor would minimize this error.
\end{itemize}

\begin{block}{Theorem: Expectation and MSE}
  If $X$ is a random variable, then the value $b$ that minimizes $E\big[(X - b)^2\big]$ is $b = E[X]$:
  $$
  \argmin_b E\big[(X - b)^2\big] = E[X].
  $$
  
  \mode<article>{
  \begin{proof}
    \begin{align*}
      E\big[(X - b)^2\big] &= E\big[(X - E(X) + E(X) - b)^2\big] \\
      &= E\big[(X - E[X]) + (E[X] - b)^2\big] \\
      &= E\big[(X - E[X])^2\big] + \big(E[X] - b\big)^2 + 2E\big[(X - E[X])(E[X] - b)\big].
    \end{align*}
    \begin{itemize}
      \item Above, we have just expanded the square, and used the fact that both $E[X]$ and $b$ are constants.
      \item Using this same idea, because $(E[X] - b)$ is just a constant, we can pull it out of the expectation in the last term, leaving $c \times E\big[X - E[X]\big]$ for some constant $c$, but $E\big[X - E[X]\big] = 0$.
      \item Thus, we have the equality
      $$
      E\big[(X - b)^2\big] = E\big[(X - E[X])^2\big] + \big(E[X]-b\big)^2.
      $$
      \item As a reminder, we are trying to minimize this quantity with respect to $b$.
      \item The first term does not involve $b$, so it has no role in the minimization.
      \item The second term is a quadratic function of $b$, meaning it has a unique global minimum.
      \item You could take the derivative and set equal to zero at this point, or just notice that if $b = E[X]$, the quadratic function is equal to $0$, which is the smallest it could be.
      \item Thus, $b = E[X]$ minimizes the MSE.
    \end{itemize}
  \end{proof}
  }
  
\end{block}

\mode<presentation>{
  \emph{Proof:}
}

\end{frame}


\begin{frame}[allowframebreaks]{Some comments on expected values}

\begin{itemize}
  \item An important thing to notice about the theorem for linear combinations is that we do not require independence.
  \item The last example demonstrates this principle. Though $I_n$ is Bernoulli distributed, $\sum_n I_n$ is \alert{NOT} binomial distributed, because the $I_n$ are not independent.
  \item As an example, if our word is $TATA$, then $I_1 = 1$ implies that $I_2 = 0$, since a $TATA$ at position 1 implies that the second letter starts with $A$, and thus $TATA$ cannot occur at position $2$.
  \item Despite this, we can still calculate the expected value of a sum by taking the sum of expected values.

\framebreak

    \item The expected value can be used as an indication of the central value of the density or frequency function.
    \item Because of this, the expected value is sometimes referred to as a \alert{location parameter}.
    \item The expected value is not the only type of location parameter. For instance, the \emph{median} is also a type of location parameter.
    \item We have seen a lot of parallel between the expected value of a discrete random variable and that of a continuous random variable. This is not a coincidence. 
    \item Specifically, we generally just ``swap" and integration with summation, and pdf with pmfs.
    \item With a more rigorous definition of expectation, we could define expectation as a \alert{Lebesgue-Stieltjes} integral, with respect to some measure $P$.
    \item That is, $E(X) = \int_\Omega XdP$, where $P$ is a probability measure. If the probability measure is a counting measure, then the integral \emph{is} a sum. 
    \item Note that this definition does not require the existence of a pdf; in fact, there distributions where the expectation is well-defined, but the pdf is not. These types of distributions do not come up often in standard examples. 
  \end{itemize}
\end{frame}

\section{Variance and Standard Deviation}

\begin{frame}[allowframebreaks]{Variance}
  \begin{itemize}
    \item The expected value is useful for summarizing the average or expected behavior of a random variable.
    \item We are also often interested in the ``spread" of a random variable.
    \item That is, if the expected value is the center (or location) of a distribution, we want an indication of how dispersed a distribution is around this center.
    \item The two most common ways to express this idea is the \alert{variance} and \alert{standard deviation} of a random variable.
  \end{itemize}
  
  \framebreak
  
  \begin{block}{Definition: Variance}
    If $X$ is a random variable with expected value $E(X)$, then the \alert{variance} of $X$ is
    $$
    \Var(X) = E\Big[\big(X - E(X)\big)^2\Big],
    $$
    provided the expectation exists.
  \end{block}
  
  \framebreak
  
  \begin{itemize}
    
    \item Letting $\mu = E[X]$, we can use the identity $g(x) = (X - \mu)^2$, and our expression for $E\big[g(X)\big]$ to get a way of calculating the variance.
  
    \item If $X$ is a discrete random variable, then by Theorem~4.1, 
    $$
    \Var(X) = \sum_{i}(x_i - \mu)^2\,p(x_i),
    $$
    \item If $X$ is a continuous random variable, then
    $$
    \Var(X) = \int_{-\infty}^\infty (x- \mu)^2 f(x)\, dx
    $$
  \end{itemize}
  
  \framebreak
  
  \begin{block}{Definition: Standard deviation}
    If $X$ is a random variable, then the standard deviation of $X$ is the square-root of the variance, provided it exists. 
  \end{block}
  
  \begin{itemize}
    \item The variance is often denoted by $\sigma^2$, and the standard deviation $\sigma$.
    \item Because $\big(X - E(X)\big)^2 \geq 0$, $\Var(X) \geq 0$.
    \item Formally, the variance is the mean of the squared distance between $X$ and $E[X]$. If most values of $X$ are close to the mean, this value is small; and vice-versa if most values of $X$ are far away from $E[X]$.
    \item By this definition, the units for the variance are squared units. 
    \item That is, if $X$ is measured in meters, then the variance is measured in square-meters, and the standard deviation is measured in meters. 
  \end{itemize}
  
  \framebreak
  
  \begin{block}{Theorem~4.4: linear transformation of a single variable}
    Let $X$ be a random variable, and assume that $\Var(X)$ exists. Then if $Y = a + bX$, then $\Var(X) = b^2\Var(X)$.
    
    \mode<article>{
    \begin{proof}
      \begin{align*}
        E\Big[\big(Y - E(Y)\big)^2\Big] &= E\Big[\big(a + bX - (a + bE[X]\big)^2\Big] \\
                                        &= E\Big[b^2\big(X - E[X]\big)^2\Big] \\
                                        &= b^2E\Big[\big(X - E[X]\big)^2\Big] \\
                                        &= b^2 \Var(X)
      \end{align*}
    \end{proof}
    }
  \end{block}
  
  \mode<presentation>{\emph{Proof.}}
  
  \framebreak
  
  \begin{itemize}
    \item This result makes a lot of sense: adding a constant only ``shifts" a distribution, it does not affect the spread.
    \item The multiplier does change the spread, and because we're squaring the difference, the multiplier is also squared.
    \item From this result, we can also see that the standard deviation also changes in a natural way. 
    \item Specifically, if $\sigma_Y, \sigma_X$ denote the standard deviations of $X$ and $Y$, respectively, then
    $$
    \sigma_Y = |b|\sigma_X.
    $$
    \item We take the absolute value, because variance and standard deviation are always positive, though the multiplier $b$ might be negative.
  \end{itemize}
  
  \framebreak
  
  \begin{exampleblock}{Example: Bernoulli distribution}
    Let $X$ be a Bernoulli$(p)$ distributed random variable. What is the variance of $X$?
    
    \mode<article>{\emph{Solution.} We'll calculate the variance using the definition of expectation. For a discrete random variable, that means summing the possible values by the corresponding probabilities:
    \begin{align*}
      \Var(X) &= \sum_{x} (x-p)^2 \, p(x) \\
              &= (0 - p)^2 (1-p) + (1 - p)^2 (p) \\
              &= p^2(1-p) + p(1-p)^2 \\
              &= p^2 - p^3 + p - 2p^2 + p^3 \\
              &= p(1 - p).
    \end{align*}
    \begin{itemize}
      \item Note that $p(1-p)$ is a quadratic function of $p$, that is maximized at $p = 1/2$.
      \item When $p = 0$ or $p = 1$, the variance is $0$, because the value will have value $X = 0$ or $X = 1$ with probability $1$.
    \end{itemize}
    }
  \end{exampleblock}
  
  \framebreak
  
  \begin{exampleblock}{Example: Normal distribution}
    Let $X \sim N(\mu, \sigma^2)$. What is $\Var(X)$? 
    
    \mode<article>{
      \emph{Solution.} Since we are already familiar with the normal distribution, we suspect that the variance is $\sigma^2$, both because it's the standard notation for variance, and because we call $\sigma$ the standard deviation parameter for the distribution. However, let's quickly demonstrate that this is indeed the case.
      If we denote $E(X) = \mu$, then
      \begin{align*}
        \Var(X) &= E\Big[\big(X - E[X]\big)^2\Big] \\
                &= \frac{1}{\sigma\sqrt{2\pi}} \int_{-\infty}^\infty (x - \mu)^2 e^{-\frac{1}{2}\big(\frac{x - \mu}{\sigma}\big)^2} \, dx.
      \end{align*}
      Making the change of variables $z = (x - \mu) / \sigma$, we have 
      \begin{align*}
        \Var(X) &= \frac{\sigma^3}{\sigma\sqrt{2\pi}} \int_{-\infty}^\infty z^2 e^{-z^2/2} \, dz \\
                &= \frac{2\sigma^2}{\sqrt{2\pi}} \int_0^\infty z^2 e^{-z^2/2} \, dz.
      \end{align*}
      We now make another transformation, $u = z^2 / 2$, which implies that $du = zdu$ and $z = \sqrt{2u}$, and thus
      \begin{align*}
        \Var(X) &= \frac{2\sigma^2}{\sqrt{2\pi}} \int_{0}^\infty \frac{2u}{\sqrt{2u}}e^{-u} \, du \\
                &= \frac{2\sigma^2}{\sqrt{\pi}} \int_0^\infty u^{1/2} e^{-u} \, du
      \end{align*}
      Now we can use the Gamma-distribution to help us evaluate this integral. Namely we need to find values of $\alpha$ and $\lambda$ such that the integrand matches the pdf of a Gamma$(\alpha, \lambda)$ distribution:
      $$
      f(u) = \frac{\lambda^\alpha}{\Gamma(\alpha)} u^{\alpha - 1}e^{-\lambda u}, \quad u \geq 0.
      $$
      Thus if we let $\alpha = 3/2$ and $\lambda = 1$, then we find that the integral on the right hand side is
      \begin{align*}
        \Var(X) &= \frac{2\sigma^2}{\sqrt{\pi}} \times \frac{\Gamma(\alpha)}{\lambda^\alpha} \\
                &= \frac{2\sigma^2}{\sqrt{\pi}} \times \Gamma(3/2) \\
                &= \frac{2\sigma^2}{\sqrt{\pi}} \frac{\sqrt{\pi}}{2} = \sigma^2.
      \end{align*}
    }
    
  \end{exampleblock}
  
  \framebreak
  
  \begin{itemize}
    \item Using the definition of variance, we will derive a very famous inequality.
  \end{itemize}
  
  \begin{block}{Theorem~4.5: Chebyshev's Inequality}
    Let $X$ be a random variable with $E[X] = \mu$, and $\Var(X) = \sigma^2$. Then for any $t > 0$,
    $$
    P\big(|X - \mu| > t\big) \leq \frac{\sigma^2}{t^2}.
    $$
    \mode<article>{
    \begin{proof}
      The proof of the inequality is rather trivial using Markov's inequality. Let $Y = (X - \mu)^2$. Because $Y$ is non-negative, we can use the most standard version of Markov's inequality:
      \begin{align*}
        P(Y > t^2) &\leq \frac{E[Y]}{t^2} \\
        P\big((X - \mu)^2 > t^2\big) &\leq \frac{E\big[(X - \mu)^2\big]}{t^2}\\
        P\big(|X - \mu| > t \big) &\leq \frac{\sigma^2}{t^2}.
      \end{align*}
    \end{proof}
    }
  \end{block}
  
  \framebreak
  
  \begin{itemize}
    \item This theorem bounds the probability that the difference between $X$ and $E[X]$ is larger than $t$.
    \item If $\sigma^2$ is small, then the probability that $X$ deviates far away from the mean is also small.
    \item By letting $t = k\sigma$, we get a bound on the probability that a variable will be $k$-standard deviations away from the mean:
    $$
    P\big(|X - \mu| > k\sigma\big)\leq \frac{1}{k^2},
    $$
    \item For instance, the probability that any arbitrary random variable $X$ will be more than $4\sigma$ away from $E[X]$ is less than $1/16$.
    \item While applicable to all random variables with well-defined variances, it is not the most optimal bound we can achieve.
    \item For instance, if $X \sim N(\mu, \sigma^2)$, then $P(|X - \mu| > 1.96\times \sigma) = 0.05 < 1/4$
  \end{itemize}
  
  \begin{block}{Corollary: zero variance}
    Let $X$ be a random variable with $\Var(X) = 0$. Then $P(X = \mu) = 1$.
    
    \mode<article>{
    \begin{proof}
      Suppose that $P(X = \mu) \neq 1$. Since $P$ is a probability measure, we can deduce $P(X = \mu) < 1$ from this assumption. Thus, there must exist some $\epsilon > 0$ such that $P(|X - \mu| > \epsilon) > 0$. However, this leads us to a contradiction: using Chebyshev's inequality, we know that for all $\epsilon > 0$, $P(|X - \mu| > \epsilon) = 0$, and therefore our assumption must be false, implying that $P(X = \mu) = 1$.
    \end{proof}
    }
    
  \end{block}
  
  \framebreak
  
  \begin{block}{Theorem~4.6: Variance Calculation}
    Let $X$ be a random variable such that $\Var(X)$ exists. Then
    $$
    \Var(X) = E(X^2) - \big[E(X)\big]^2 = E(X^2) - \mu^2,
    $$
    where $\mu = E(X)$.
  \end{block}
  
  \mode<article>{
  \begin{proof}
    \begin{align*}
      \Var(X) &= E\big[(X - \mu)^2\big] \\
              &= E\big[X^2 - 2\mu X + \mu^2\big] \\
              &= E(X^2) - 2\mu E(X) + \mu^2 \\
              &= E(X^2) - 2\mu^2 + \mu^2 \\
              &= E(X^2) - \mu^2.
    \end{align*}
  \end{proof}
  }
  
  \framebreak
  
  \begin{itemize}
    \item Theorem~4.6 is sometimes useful to help us calculate the variance of a random variable. 
    \item Other times, the variance is known, and the theorem helps us calculate $E(X^2)$. 
  \end{itemize}
  
  \begin{exampleblock}{Example: Uniform distribution}
    Let $X \sim U(0, 1)$. Use Theorem~4.6 to find $\Var(X)$.
    
    \mode<article>{
    \emph{Solution.} the pdf of $X$ is $f(x) = 1[0 \leq x \leq 1]$. Then
    $$
    E[X] = \int_{-\infty}^\infty x f(x) \,dx = \int_0^1 x \, dx = 1/2.
    $$
    $$
    E[X^2] = \int_{-\infty}^\infty x^2 f(x) \,dx = \int_0^1 x^2 \, dx = 1/3.
    $$
    Thus, 
    $$
    \Var(X) = E[X^2] - \big(E[X]\big)^2 = \frac{1}{3} - \left(\frac{1}{2}\right)^2 = \frac{1}{12}.
    $$
    }
  \end{exampleblock}
  
\end{frame}

\section{Bias-Variance Tradeoff}

\begin{frame}[allowframebreaks]{Measurement Error}

\begin{itemize}
  \item Often, values of interest cannot be known precisely, but instead must be determined by experimental procedures.
  \item For instance: measurements of weight, length, voltage, or intervals of time can be complex, and generally involve potential sources of error.
  \item The National Institute of Standards and Technology (NIST) in the US are charged with developing and maintaining measurement standards.
  \item Statisticians have historically been employed by these organizations to help with this endeavor.
  
  \framebreak 
  
  \item Typically, there are two main types of measurement error: \alert{random} vs \alert{systematic}.
  \item For instance, a sequence of repeated independent measurements made from the same instrument or experimental procedure may not give the same value each time. These uncontrollable differences are modeled as \alert{random} error. 
  \item However, there may be a \alert{systematic} error that affects all measurements, such as poorly calibrated instruments, or errors that are associated with the method of measurement.
  \item Suppose that the true value of a quantity being measured is $x_0$. We have a random measurement $X$, which is modeled as 
  $$
  X = x_0 + \beta + \epsilon.
  $$
  \item Here, $\beta$ is the systematic error, and $\epsilon$ is the random component of the error.
\end{itemize}
  \framebreak
  
  \begin{block}{Definition: Bias}
    Let $x_0$ be the true value of a measurement, modeled as a random variable $X$ such that
    $$
    X = x_0 + \beta + \epsilon,
    $$
    where $E(\epsilon) = 0$, $\Var(\epsilon) = \sigma^2$. Then, we have
    $$
    E[X] = x_0 + \beta.
    $$
    The value $\beta = E(X - x_0)$ is called the \alert{bias} of the random variable, and we say that $X$ is an unbiased estimate of $x_0$ if $\beta = 0$. 
  \end{block}
  
  \framebreak
  
\begin{itemize}
  \item The two factors that impact the quality of our estimator is the bias $\beta$ and the variance $\sigma^2$. 
  \item If both $\beta = 0$ and $\sigma^2 = 0$, then we get a perfect measurement.
  \item Ideally, we want an estimator that minimizes the bias and the variance, though as we will see (Math 4451) there is a principle known as the \alert{bias-variance} trade-off, which suggests that efforts to minimize bias often result in larger variance (and vice-versa).
  \item Many approaches in statistics we will cover next semester aim at finding estimators that are unbiased ($\beta = 0$), while having minimum variance as possible (that is, the minimum-variance unbiased estimator (MVUE)).
\end{itemize}

\framebreak

\begin{block}{Theorem~4.7: Mean Squared Error}
  Let $X$ be a random variable representing a random estimate for value $x_0$. The mean-squared error of the estimator $X$ is defined as $\text{MSE}(X) = E\big[(X-x_0)^2\big]$. If $\beta$ is the bias of the estimator and $\sigma^2$ the variance, then
  $$
  \text{MSE}(X) = \beta^2 + \sigma^2.
  $$
  \mode<article>{
  \begin{proof}
    \begin{align*}
      E\big[(X-x_0)^2\big] &= \Var(X - x_0) + \big[E(X - x_0)\big]^2 \\
      &= \Var(X) + \beta^2 \\
      &= \sigma^2 + \beta^2.
    \end{align*}
  \end{proof}
  }
\end{block}

\end{frame}

% \begin{frame}[allowframebreaks]{}
% 
% \end{frame}

\newcommand\acknowledgments{
\begin{itemize}
\item   Compiled on {\today} using \Rlanguage version 4.5.1.
\item   \parbox[t]{0.75\textwidth}{Licensed under the \link{http://creativecommons.org/licenses/by-nc/4.0/}{Creative Commons Attribution-NonCommercial license}.
    Please share and remix non-commercially, mentioning its origin.}
    \parbox[c]{1.5cm}{\includegraphics[height=12pt]{../cc-by-nc}}
\item We acknowledge \link{https://jeswheel.github.io/4450_f25/acknowledge.html}{students and instructors for previous versions of this course / slides}.
\end{itemize}
}

\mode<presentation>{
\begin{frame}[allowframebreaks=0.8]{References and Acknowledgements}
  
\bibliography{../bib4450}

\vspace{3mm}

\acknowledgments

\end{frame}
}

\mode<article>{

\newpage

{\bf \Large \noindent Acknowledgments}

\acknowledgments

\newpage

\bibliography{../bib4450}

}



\end{document}







